\chapter{The Asset Managers' Strategies}\label{ch:s1:the-asset-managers-strategies}

An index is a paper portfolio: it trades for free, never holds cash and rebalances at closing prices. A fund that tracks it pays for every one of those
things, and Frino and Gallagher, studying S\&P 500 \omterm{def:s1:index-rebalancing:fund}{index funds}, found tracking error unavoidable because of market frictions; they also found that the
\omterm{def:s1:index-rebalancing:fund}{index funds}, on average, beat active funds after expenses. The craft of the largest asset managers is keeping a portfolio close to a benchmark cheaply,
tilting it on purpose, and moving money between portfolios without losing it on the way. On this chapter's synthetic market, holding the largest 100
of about a thousand names with optimised weights tracks the index to 1.64\% a year, the largest 400 to 0.32\%; a transition of 59\% of a \$2 billion fund
costs 27.6 basis points if done in a day and least in cost plus risk if spread over two or three. The build is \texttt{firm.assetmgr}.

\section{Indexing: replication and sampling}

\begin{definition}[Full replication, sampled replication]\label{def:s1:the-asset-managers-strategies:replication}
\emph{Full replication}\index{full replication} holds every constituent of an index at its index weight. \emph{Sampled replication}\index{sampled
replication} holds a subset of the constituents, with weights chosen so that the portfolio's risk exposures match the index's and its forecast tracking
error is as small as possible, trading lower costs for some tracking error.
\end{definition}

A cap-weighted index of \texttt{firm.synthmkt}'s listed names drifts with prices by itself: holding it exactly needs trades only when names enter or
leave, a one-way turnover of 6.4\% a year here. Sampling needs a risk model to decide which names stand for which. The chapter uses the synthetic
market's own truth, the best case: each name's beta on the market factor (with its conditional variance), industry factors at 8\% a year, the three
styles at their planted volatilities with point-in-time exposures, and each name's specific volatility (Book~7, chapter~24 estimates such a model from
returns). \Cref{lst:s1:the-asset-managers-strategies:sampled} keeps the $n$ largest names and solves a quadratic programme for the long-only weights
that minimise the ex-ante tracking error; the book is rebalanced every 63 days and drifts with returns in between.

\begin{center}
\small
\begin{tabular}{@{}lccccc@{}}
\toprule
names held & 25 & 50 & 100 & 200 & 400\\
\midrule
realised tracking error & 4.84\% & 2.96\% & 1.64\% & 0.81\% & 0.32\%\\
ex-ante tracking error (mean) & 4.70\% & 2.81\% & 1.52\% & 0.71\% & 0.26\%\\
one-way turnover a year & 1.34 & 1.13 & 0.92 & 0.58 & 0.31\\
cost a year at 10 bp per unit traded (bp) & 13.4 & 11.3 & 9.2 & 5.8 & 3.1\\
\bottomrule
\end{tabular}
\end{center}

Tracking error falls fast with the number of names (\cref{fig:s1:the-asset-managers-strategies:sampling}), by 40\% to 60\% for each doubling, because
each added name is smaller and the specific risk left out falls with the index weight left out. The ex-ante figure understates the realised one a little
even with the true model: between rebalances the weights drift and listings change. Costs go the other way from what one might expect: the small samples
trade more, because their weights must be re-optimised as the index moves. Over the ten years the 25-name sample happened to beat the index by 3.0\% a
year; with a tracking error of 4.84\%, that is noise, and a sample that holds only the largest names is a bet on large caps.

\begin{omfigure}
\begin{tikzpicture}
\begin{axis}[width=10cm,height=5.2cm,xmode=log,ymode=log,xlabel={\small names held},ylabel={\small tracking error a year (\%)},
  tick label style={font=\footnotesize},xmin=20,xmax=500,ymin=0.2,ymax=6,grid=major,grid style={gray!25},legend pos=north east,
  legend style={font=\scriptsize},xtick={25,50,100,200,400},xticklabels={25,50,100,200,400},ytick={0.25,0.5,1,2,4},
  yticklabels={0.25,0.5,1,2,4},table/col sep=comma]
\addplot[omThm,thick,mark=*] table[x=names,y=te_pct]{figdata/strategies-1/29-the-asset-managers-strategies/sampling.csv};
\addplot[omDef,thick,dashed,mark=square*] table[x=names,y=te_ex_pct]{figdata/strategies-1/29-the-asset-managers-strategies/sampling.csv};
\legend{realised,ex ante}
\end{axis}
\end{tikzpicture}

\omcaption{\omterm{def:s1:the-asset-managers-strategies:replication}{Sampled replication} of the synthetic cap-weighted index: realised and ex-ante tracking error against the number of largest names held, with
long-only weights minimising tracking error under the market's true risk model, rebalanced every 63 days (log scales). Data:
\texttt{s1\_assetmgr.replicate}.}\label{fig:s1:the-asset-managers-strategies:sampling}
\end{omfigure}

Real \omterm{def:s1:index-rebalancing:fund}{index funds} add what the model leaves out: dividends to reinvest, cash flows from investors, index changes announced in advance (chapter~10), and
the choice of when to trade them. Frino, Gallagher and Oetomo found that passive funds benefit from less rigid rebalancing, that enhanced \omterm{def:s1:index-rebalancing:fund}{index funds}
rebalance earlier and more patiently around index revisions and pay less for it, and that passive funds, when they deviate from the index, tend to
overweight larger, more liquid stocks.

\section{Factor products}

\begin{definition}[Factor product]\label{def:s1:the-asset-managers-strategies:factor}
A \emph{factor product}\index{factor product} is a rules-based portfolio, usually long-only, that tilts a broad index toward a style (value, momentum,
quality, low volatility, size) by a published rule, sold as a fund or an index at a fee between an \omterm{def:s1:index-rebalancing:fund}{index fund}'s and an active manager's.
\end{definition}

A \omterm{def:s1:the-asset-managers-strategies:factor}{factor product} sells exposure, not return. The chapter tilts the index by $\exp(kz)$ on each name's point-in-time book-to-price z-score:

\begin{center}
\small
\begin{tabular}{@{}lccc@{}}
\toprule
tilt $k$ & 0.25 & 0.5 & 1.0\\
\midrule
value exposure (z-score units) & 0.26 & 0.52 & 1.03\\
tracking error & 1.57\% & 3.08\% & 6.00\%\\
excess return a year & 0.07\% & 0.07\% & 0.06\%\\
\bottomrule
\end{tabular}
\end{center}

The exposure is delivered as specified, with a tracking error proportional to the tilt; the return is not. The synthetic market plants a value premium,
but a book-to-price measured with today's price is short the planted momentum (chapter~6 found the same), and the two cancel: the product earns almost
nothing over ten years. A \omterm{def:s1:the-asset-managers-strategies:factor}{factor product}'s buyer needs a view on the factor, not on the product.

\section{Benchmark-relative management}

Active managers measured against a benchmark live in the space these tools define. Their risk is tracking error, their skill is information ratio
(Book~1, chapter~3). The mean--variance answer to how much tracking error to take is simple: a manager with information ratio IR, penalising active
variance with a coefficient $\lambda$, maximises $\text{IR}\,\omega - \lambda\omega^2$ and takes $\omega = \text{IR}/(2\lambda)$; Grinold and Kahn's
\emph{Active Portfolio Management} develops the whole framework. The sampled index is the zero-skill end of that line, the \omterm{def:s1:the-asset-managers-strategies:factor}{factor product} a
constant tilt, and the active manager a time-varying one. The chapter's tools apply to all three: the same risk model sets tracking error, the same
optimiser keeps exposures inside limits, and the same costs decide how often to rebalance.

\section{Transition management}

\begin{definition}[Transition management]\label{def:s1:the-asset-managers-strategies:transition}
\emph{Transition management}\index{transition management} is the moving of a portfolio from one set of holdings to another (a new manager, a new
benchmark, a new allocation) at the lowest implementation shortfall, balancing the market impact of trading fast against the risk of holding the old
portfolio while the new one is built.
\end{definition}

A pension fund replacing its 100-name index sample with the value product (a tilt of 0.5) must trade 59\% of a \$2 billion fund, one way. Trading it in
one day pays the half spread (5 basis points, assumed) and a square-root impact on each name's daily volatility (\texttt{firm.tcost}, with $\eta = 0.5$,
assumed); spreading it over more days cuts the impact but leaves the unexecuted part exposed to the active risk between the two portfolios, 3.46\% a
year (21.8 basis points a day). \Cref{lst:s1:the-asset-managers-strategies:sampled} includes the planner.

\begin{center}
\small
\begin{tabular}{@{}lccccc@{}}
\toprule
days to complete & 1 & 2 & 3 & 5 & 10\\
\midrule
expected cost (bp of the fund) & 27.6 & 21.3 & 18.5 & 15.6 & 12.8\\
standard deviation of the outcome (bp) & 21.8 & 24.4 & 27.2 & 32.3 & 42.8\\
cost plus one standard deviation (bp) & 49.4 & 45.6 & 45.6 & 48.0 & 55.6\\
\bottomrule
\end{tabular}
\end{center}

The trade-off (\cref{fig:s1:the-asset-managers-strategies:transition}) has a broad minimum at two to three days for a planner who weighs one standard
deviation of risk as heavily as a basis point of cost; a planner who cares only about expected cost would take ten days and a 43 basis point standard
deviation. Real transitions reduce both by crossing: the old and new portfolios often share names, and trades between them, or against other clients'
flows, need not touch the market (Book~7, chapter~27).

\begin{omfigure}
\begin{tikzpicture}
\begin{axis}[width=10cm,height=5.2cm,xlabel={\small days to complete the transition},ylabel={\small basis points of the fund},
  tick label style={font=\footnotesize},xmin=0,xmax=11,ymin=0,ymax=60,grid=major,grid style={gray!25},legend columns=3,
  legend style={font=\scriptsize,at={(0.5,-0.3)},anchor=north,draw=none,fill=none,/tikz/every even column/.append style={column sep=8pt}},
  table/col sep=comma]
\addplot[omThm,thick,mark=*] table[x=days,y=cost_bp]{figdata/strategies-1/29-the-asset-managers-strategies/transition.csv};
\addplot[omDef,thick,mark=square*] table[x=days,y=risk_bp]{figdata/strategies-1/29-the-asset-managers-strategies/transition.csv};
\addplot[black!60,thick,dashed,mark=triangle*] table[x=days,y=total_bp]{figdata/strategies-1/29-the-asset-managers-strategies/transition.csv};
\legend{expected cost,standard deviation,cost plus one standard deviation}
\end{axis}
\end{tikzpicture}

\omcaption{A transition of 59\% of a \$2 billion synthetic fund from a 100-name index sample to a value tilt: expected cost (half spread plus square-root
impact, assumed parameters) and the standard deviation from holding the unexecuted part, by the number of days taken. Data:
\texttt{s1\_assetmgr.transition\_plan}.}\label{fig:s1:the-asset-managers-strategies:transition}
\end{omfigure}

\section{Strategy files}

\begin{strategyfile}{Sampled index replication}\label{strat:s1:the-asset-managers-strategies:index}
\sfield{Who pays you, and why} Investors paying a small fee for the market's return at low cost and low tracking error.
\sfield{Instruments and venues} Index constituents; futures for cash equitisation.
\sfield{Signal} None: the index and a risk model.
\sfield{Sizing and execution} \omterm{def:s1:the-asset-managers-strategies:replication}{Full replication} where costs allow, sampling otherwise; trades at the close on index changes, or patiently around them.
\sfield{Costs} Turnover from index changes and re-optimisation.
\sfield{How it dies} It does not; fees compete toward zero.
\sfield{Horizon, capacity, infrastructure} Permanent; very large capacity; index data, a risk model, an optimiser.
\sfield{Backtest honestly} The index as it was, with its changes and their announcement dates; dividends and cash.
\sfield{Sources} Frino and Gallagher (2001); this chapter: 1.64\% tracking error with 100 names, 0.32\% with 400.
\end{strategyfile}

\begin{strategyfile}{Enhanced indexing and factor products}\label{strat:s1:the-asset-managers-strategies:enhanced}
\sfield{Who pays you, and why} Investors wanting a factor premium or small excess returns at a known tracking error.
\sfield{Instruments and venues} Index constituents.
\sfield{Signal} A published tilt rule, or patient trading around index changes.
\sfield{Sizing and execution} Tilts sized to a tracking-error budget; rebalanced on schedule.
\sfield{Costs} Moderate turnover.
\sfield{How it dies} Factors that stop paying; crowding.
\sfield{Horizon, capacity, infrastructure} Years; large capacity.
\sfield{Backtest honestly} Point-in-time exposures; the factor's interaction with others (value against momentum).
\sfield{Sources} Frino, Gallagher and Oetomo (2005); this chapter: exposure delivered, return near zero.
\end{strategyfile}

\begin{strategyfile}{Transition management}\label{strat:s1:the-asset-managers-strategies:transition}
\sfield{Who pays you, and why} Asset owners changing managers or allocations, who pay to avoid the shortfall of doing it badly.
\sfield{Instruments and venues} Both portfolios' holdings; crossing networks; futures for interim exposure.
\sfield{Signal} The trade list and the active risk between the two portfolios.
\sfield{Sizing and execution} A schedule balancing impact against risk; crossing first.
\sfield{Costs} Spread, impact and the risk of the unexecuted part.
\sfield{How it dies} It does not; it competes on execution.
\sfield{Horizon, capacity, infrastructure} Days; a cost model and a risk model.
\sfield{Backtest honestly} Implementation shortfall against the prices at the decision, not at completion.
\sfield{Sources} This chapter: 27.6 basis points in one day, least cost plus risk over two or three days.
\end{strategyfile}

\section{Tutorial: invisible craft}\label{tut:s1:the-asset-managers-strategies:tut}

\begin{tutorial}
\textbf{Goal.} Replicate the synthetic index by sampling, measure tracking error and costs by the number of names, build a value product, and plan a
transition. \textbf{End state:} the tables and the two figures.
\begin{tutsteps}
\item \textbf{Sampling, drift, tilts and transitions}.
\omcode{code/firm/assetmgr/firm_assetmgr.py}{44}{80}{Sampled replication, drift, tilts and the transition planner.}\label{lst:s1:the-asset-managers-strategies:sampled}
\item \textbf{The replication backtest}: optimise every 63 days, drift in between.
\omcode{code/strategies-1/29-the-asset-managers-strategies/python/s1_assetmgr.py}{65}{80}{Sampled replication over ten years.}\label{lst:s1:the-asset-managers-strategies:replicate}
\item \textbf{Run} \texttt{replicate(n)} for five sizes, \texttt{full\_turnover()}, \texttt{value\_product(k)} and \texttt{transition\_plan()}, and
\texttt{fig\_assetmgr.py}.
\end{tutsteps}
\textbf{What to change next.} Estimate the risk model from returns (Book~7, chapter~24) instead of using the truth and compare tracking errors; select
names by stratifying on industry instead of size; cross the transition's overlapping names before trading.
\end{tutorial}

\section{Build: the asset manager's toolkit}\label{bld:s1:the-asset-managers-strategies:assetmgr}

\begin{build}
\bfield{Purpose} Index weights, \omterm{def:s1:the-asset-managers-strategies:replication}{sampled replication}, drift of held weights, factor tilts and a transition planner.
\bfield{Interface} \texttt{cap\_weights(price, shares, listed)}, \texttt{covariance(B, F, D)}, \texttt{sampled(b, Sigma, n)}, \texttt{drift\_active(R, w0,
b\_path, start, end)}, \texttt{tilt(b, z, k)}, \texttt{transition(value, adv, sigma, half\_spread, eta, active\_sd, days)}.
\bfield{Rules} Long-only weights summing to one; the index from the previous close; costs from firm.tcost.
\bfield{Acceptance tests} \texttt{code/firm/assetmgr/tests/}: cap weights and a covariance by hand, sampling with all names giving zero tracking error and
with ten names a positive one, drift, a neutral tilt, and a transition's cost falling and risk rising with days.
\bfield{Stretch} Stratified sampling; an estimated risk model; crossing in transitions.
\end{build}

\begin{omsources}
\item A.~Frino and D.~R.~Gallagher, ``Tracking S\&P 500 index funds'', \emph{Journal of Portfolio Management} 28(1), 2001.
\item A.~Frino, D.~R.~Gallagher and T.~N.~Oetomo, ``The index tracking strategies of passive and enhanced index equity funds'', \emph{Australian Journal
of Management} 30(1), 2005.
\item R.~C.~Grinold and R.~N.~Kahn, \emph{Active Portfolio Management}, 2nd ed., McGraw-Hill, 2000.
\end{omsources}

\section{Exercises}

\begin{exercise}[$\star$]\label{exo:s1:the-asset-managers-strategies:1}
If tracking error fell with the square root of the number of names, by what factor would doubling the names cut it? What does the table show from 25 to 50
names, and from 100 to 200?
\end{exercise}

\begin{exercise}[$\star$]\label{exo:s1:the-asset-managers-strategies:2}
Holding the index exactly turns over 6.4\% a year. At 10 basis points per unit traded, what does it cost?
\end{exercise}

\begin{exercise}[$\star$]\label{exo:s1:the-asset-managers-strategies:3}
Why does a \omterm{def:s1:the-asset-managers-strategies:factor}{factor product}'s tracking error grow with the tilt while its return need not?
\end{exercise}

\begin{exercise}[$\star\star$]\label{exo:s1:the-asset-managers-strategies:4}
The transition's active risk is 3.46\% a year. What is it per day, and why is that the standard deviation of a one-day transition?
\end{exercise}

\begin{exercise}[$\star\star$]\label{exo:s1:the-asset-managers-strategies:5}
Why is the realised tracking error of a sample larger than its ex-ante tracking error even with the true risk model?
\end{exercise}

\begin{exercise}[$\star\star$]\label{exo:s1:the-asset-managers-strategies:6}
With a value premium of 3\% a year per unit of exposure, what would a tilt with an exposure of 0.52 earn if value were uncorrelated with everything else?
Why did it earn 0.07\%?
\end{exercise}

\begin{exercise}[$\star\star\star$]\label{exo:s1:the-asset-managers-strategies:7}
\emph{Coding.} Change the transition's objective to cost plus two standard deviations. How many days does the planner choose now?
\end{exercise}

\begin{exercise}[$\star\star\star$]\label{exo:s1:the-asset-managers-strategies:8}
\emph{Find the flaw.} ``Our 50-stock index sample beat the index by 2\% a year for ten years, so sampling adds value.''
\end{exercise}

\section{Problem: Invisible Craft}

\begin{problem}[{Weekend problem --- an index, a tilt and a transition}]\label{pb:s1:the-asset-managers-strategies:1}
The chapter's synthetic market and the public record.

\medskip\noindent\textbf{Part I --- Indexing.}
\begin{enumerate}
\item Define full and \omterm{def:s1:the-asset-managers-strategies:replication}{sampled replication}.
\item What did Frino and Gallagher find about \omterm{def:s1:index-rebalancing:fund}{index funds}?
\item Why does holding the index exactly need so little trading?
\item Describe the chapter's risk model and sampling optimiser.
\end{enumerate}

\medskip\noindent\textbf{Part II --- Sampling.}
\begin{enumerate}[resume]
\item Give the realised and ex-ante tracking errors by the number of names.
\item Why do small samples trade more?
\item What did the 25-name sample's excess return mean?
\item What did Frino, Gallagher and Oetomo find about rebalancing?
\end{enumerate}

\medskip\noindent\textbf{Part III --- Products and benchmarks.}
\begin{enumerate}[resume]
\item Define a \omterm{def:s1:the-asset-managers-strategies:factor}{factor product} and give its exposure, tracking error and return by tilt.
\item Why did the value product earn almost nothing?
\item How does Grinold and Kahn's framework budget tracking error?
\item Define \omterm{def:s1:the-asset-managers-strategies:transition}{transition management}.
\end{enumerate}

\medskip\noindent\textbf{Part IV --- The verdict.}
\begin{enumerate}[resume]
\item State the \emph{named result}: the tracking error of \omterm{def:s1:the-asset-managers-strategies:replication}{sampled replication} against the number of names held, and the transition's shortfall.
\item What limits the transition's speed, and what limits its patience?
\item How does crossing reduce a transition's cost?
\item What would an estimated risk model change?
\item How would you backtest an \omterm{def:s1:index-rebalancing:fund}{index fund} honestly?
\item Which strategy file depends most on the risk model?
\item How does this chapter relate to chapter~10?
\item In one sentence: what is the asset manager's craft?
\end{enumerate}
\end{problem}

\section{Interview questions}

\begin{interviewq}[$\star$ \iqroles{researcher}]\label{iq:s1:the-asset-managers-strategies:1}
What causes an \omterm{def:s1:index-rebalancing:fund}{index fund}'s tracking error?
\end{interviewq}

\begin{interviewq}[$\star\star$ \iqroles{researcher}]\label{iq:s1:the-asset-managers-strategies:2}
How would you choose which names to hold in a \omterm{def:s1:the-asset-managers-strategies:replication}{sampled replication} of a 3\,000-name index?
\end{interviewq}

\begin{interviewq}[$\star\star$ \iqroles{trader}]\label{iq:s1:the-asset-managers-strategies:3}
You must move a \$5 billion portfolio to a new manager's holdings. Walk through the plan.
\end{interviewq}

\begin{interviewq}[$\star\star$ \iqroles{developer}]\label{iq:s1:the-asset-managers-strategies:4}
What data and checks does an \omterm{def:s1:index-rebalancing:fund}{index fund} need around an index rebalance?
\end{interviewq}

\begin{interviewq}[$\star\star$ \iqroles{risk}]\label{iq:s1:the-asset-managers-strategies:5}
How would you monitor a \omterm{def:s1:the-asset-managers-strategies:factor}{factor product}'s risk, and what would make you worried?
\end{interviewq}

\begin{interviewq}[$\star\star\star$ \iqroles{researcher}]\label{iq:s1:the-asset-managers-strategies:6}
A trade list is executed evenly over $N$ days with expected cost $c_0 N^{-1/2}$ (square-root impact) and risk $\sigma_a \sqrt{\sum_{d=0}^{N-1}(1 -
d/N)^2}$. Approximate the risk for large $N$ and find the $N$ that minimises cost plus $\lambda$ times risk.
\end{interviewq}
